% Reconciled against public ORCID, publisher Crossref deposits, and author notices.
% Sources and exceptions: data/publication-refresh-2026-09-29.md.
\section{Publications}
{\fontsize{9}{11}\selectfont\color{muted}%
Bold reference numbers mark selected work and link back to the research programme. [DOI] opens the article.\par}
\subsection*{Journal articles / accepted and in press}
\begin{publications}
\pub{A01}{\guo{} (2026). Delivered-Energy and Operative-Temperature Trade-offs for Learned-Probability-Informed HVAC Supervision in a Weather Stress Grid. \emph{Energy and Buildings}. \textbf{Accepted 27 September 2026}.}
\pub{A02}{\guo{}, He K. (2026). Lower Peaks, Similar Growth: Isolating the Value and Limits of AI-Driven Thermal Sensation Modeling for Building--Grid Resilience Under Future Climates. \emph{Engineering Applications of Artificial Intelligence}. \textbf{Accepted 15 September 2026}.}
\pub{A03}{\guo{} (2026). Equal cooling setpoints, unequal thermal outcomes: Residential cooling service under future weather and ageing. \emph{Energy}, article 142403. \textbf{In press}. \doi{10.1016/j.energy.2026.142403}.}
\end{publications}

\subsection*{Journal articles / published}
\begin{publications}

\pub{J01}{Chang Y., \guo{} (2026). Temporal regime diagnostics of dining-space CO2 under occupancy loading: A long-horizon operational study with dense-grid seed evidence. \emph{Building and Environment}, \textbf{305}, 115257. Published online 14 September 2026. \doi{10.1016/j.buildenv.2026.115257}.}

\pub{J02}{\guo{}, Anselmo S., Ferrara M., Niu S., Chi B., Wang X. (2026). Automated urban energy assessment: From thermal flyover to AI-driven retrofit prioritization for sustainable cities. \emph{Sustainable Cities and Society}, \textbf{148}, 107589. \doi{10.1016/j.scs.2026.107589}.}

\pub{J03}{\guo{}, He K. (2026). Scenario-conditioned actual meteorological years (sAMY): A stochastic weather generator using multi-decadal observations. \emph{Energy and Buildings}, \textbf{364}, 117508. \doi{10.1016/j.enbuild.2026.117508}.}

\pub{J04}{Cen C., \guo{}, Chew L.W., Wong N.H. (2026). Experimental study on gender differences in thermal comfort and physiological responses in fan-assisted cooling environments. \emph{Energy and Buildings}, \textbf{360}, 117365. \doi{10.1016/j.enbuild.2026.117365}.}

\pub{J05}{\guo{}, He K. (2026). Input quality, not statistical complexity, determines climate-adapted weather file fidelity: A causal decomposition of degree-day errors. \emph{Energy}, \textbf{353}, 140867. \doi{10.1016/j.energy.2026.140867}.}

\pub{J06}{\guo{}, Aviv D. (2026). From seven points to probabilities: Ordinal learning for risk-aware thermal comfort prediction. \emph{Building and Environment}, \textbf{295}, 114426. \doi{10.1016/j.buildenv.2026.114426}.}

\pub{J07}{\guo{}, Pigliautile I., Chang Y., Qiao Q., Li Y. (2026). Toward smarter HVAC control: Machine learning reveals hidden drivers in thermal comfort databases. \emph{Energy and AI}, \textbf{24}, 100709. \doi{10.1016/j.egyai.2026.100709}.}

\pub{J08}{\guo{}, Chang Y., Li Y., Zhou Y., Qiao Q., Lai C.Y., Schuldenfrei E. (2026). Ventilation--energy trade-offs in retrofitted Hong Kong wet markets. \emph{Energy and Buildings}, \textbf{355}, 116918. \doi{10.1016/j.enbuild.2025.116918}.}

\pub{J09}{\guo{}, He K., Xu Y., Lei Y. (2026). A co-simulation methodology for integrating data-driven thermal sensation models with building energy control. \emph{Energy and Buildings}, \textbf{353}, 116745. \doi{10.1016/j.enbuild.2025.116745}.}

\pub{J10}{\guo{}, Sun R., Xu Y. (2025). Correcting the 120-watt assumption: Demographic-aware metabolic rates for energy savings and thermal comfort equity in buildings. \emph{Energy and Buildings}, \textbf{349}, 116525. \doi{10.1016/j.enbuild.2025.116525}.}

\pub{J11}{Chang Y., \guo{}, Li Y., Pigliautile I., Chi B. (2025). A data-driven qualitative review of thermal comfort studies: Bridging the gap between western and eastern perspectives. \emph{Renewable and Sustainable Energy Reviews}, \textbf{224}, 116020. \doi{10.1016/j.rser.2025.116020}.}

\pub{J12}{\guo{}, He K., Luo Y., Chang Y. (2025). Physics-Informed Neural Networks for robust thermal comfort prediction: Overcoming data quality limitations through physiological constraints. \emph{Building and Environment}, \textbf{285}, 113588. \doi{10.1016/j.buildenv.2025.113588}.}

\pub{J13}{Aviv D., \guo{}, Middel A., Meggers F. (2021). Evaluating radiant heat in an outdoor urban environment: Resolving spatial and temporal variations with two sensing platforms and data-driven simulation. \emph{Urban Climate}, \textbf{35}, 100745. \doi{10.1016/j.uclim.2020.100745}.}

\pub{J14}{Ma N., Aviv D., \guo{}, Braham W.W. (2021). Measuring the right factors: A review of variables and models for thermal comfort and indoor air quality. \emph{Renewable and Sustainable Energy Reviews}, \textbf{135}, 110436. \doi{10.1016/j.rser.2020.110436}.}

\pub{J15}{\guo{}, Ferrara M., Coleman J., Loyola M., Meggers F. (2020). Air temperature and mean radiant temperature data, collected and simulated across a radiantly-heated high-bay laboratory. \emph{Data in Brief}, \textbf{30}, 105192. \doi{10.1016/j.dib.2020.105192}.}

\pub{J16}{Teitelbaum E., Chen K.W., Meggers F., \guo{}, Houchois N., Pantelic J., Rysanek A. (2020). Globe thermometer free convection error potentials. \emph{Scientific Reports}, \textbf{10}(1), 2652. \doi{10.1038/s41598-020-59441-1}.}

\pub{J17}{\guo{}, Ferrara M., Coleman J., Loyola M., Meggers F. (2020). Simulation and measurement of air temperatures and mean radiant temperatures in a radiantly heated indoor space. \emph{Energy}, \textbf{193}, 116369. \doi{10.1016/j.energy.2019.116369}.}

\pub{J18}{\guo{}, Aviv D., Loyola M., Teitelbaum E., Houchois N., Meggers F. (2020). On the understanding of the mean radiant temperature within both the indoor and outdoor environment, a critical review. \emph{Renewable and Sustainable Energy Reviews}, \textbf{117}, 109207. \doi{10.1016/j.rser.2019.06.014}.}

\pub{J19}{Luo Y., \guo{}, Meggers F., Zhang L. (2019). Deep coaxial borehole heat exchanger: Analytical modeling and thermal analysis. \emph{Energy}, \textbf{185}, 1298--1313. \doi{10.1016/j.energy.2019.05.228}.}

\pub{J20}{\guo{}, Luo Y., Meggers F., Simonetti M. (2019). Human body exergy consumption models' evaluation and their sensitivities towards different environmental conditions. \emph{Energy}, \textbf{183}, 1075--1088. \doi{10.1016/j.energy.2019.05.045}.}

\pub{J21}{Luo Y., Zhang L., Bozlar M., Liu Z., \guo{}, Meggers F. (2019). Active building envelope systems toward renewable and sustainable energy. \emph{Renewable and Sustainable Energy Reviews}, \textbf{104}, 470--491. \doi{10.1016/j.rser.2019.01.005}.}

\pub{J22}{\guo{}, Teitelbaum E., Houchois N., Bozlar M., Meggers F. (2018). Revisiting the use of globe thermometers to estimate radiant temperature in studies of heating and ventilation. \emph{Energy and Buildings}, \textbf{180}, 83--94. \doi{10.1016/j.enbuild.2018.08.029}.}

\pub{J23}{Meggers F., \guo{}, Teitelbaum E., Aschwanden G., Read J., Houchois N., Pantelic J., Calabrò E. (2017). The Thermoheliodome -- ``Air conditioning'' without conditioning the air, using radiant cooling and indirect evaporation. \emph{Energy and Buildings}, \textbf{157}, 11--19. \doi{10.1016/j.enbuild.2017.06.033}.}

\pub{J24}{Meggers F., Aschwanden G., Teitelbaum E., \guo{}, Salazar L., Bruelisauer M. (2016). Urban cooling primary energy reduction potential: System losses caused by microclimates. \emph{Sustainable Cities and Society}, \textbf{27}, 315--323. \doi{10.1016/j.scs.2016.08.007}.}

\end{publications}

\subsection*{Conference papers and workshop contributions}
\begin{publications}
\pub{C01}{He K., \guo{} (2026). Data-driven cross-regional climate zone classification for building codes: A future-compatible framework. \emph{E3S Web of Conferences}, \textbf{716}, 10011. \textbf{Best Paper Award, IAQVEC 2026 (May)}. \doi{10.1051/e3sconf/202671610011}.}
\pub{C02}{He K., \guo{} (2025). A temporal features-enhanced mixture-of-experts approach for indoor temperature prediction. \emph{ICML CO-BUILD Workshop}, oral presentation. \textbf{1st Place, Smart Building Competition}.}
\pub{C03}{Niu S., \guo{}, Ferrara M., Anselmo S. (2025). Thermal-SAM: Adversarial prompt-based unsupervised building segmentation in thermal aerial imagery. \emph{ICML CO-BUILD Workshop}, poster. \textbf{3rd Place, NeurIPS Urban AI Workshop (Buildings Challenge)}.}
\pub{C04}{\guo{}, Chang Y., Hu D. (2025). ML-driven sensitivity analysis for lean HVAC: New insights from large-scale comfort data. \emph{ICML CO-BUILD Workshop}, poster.}
\pub{C05}{\guo{}, Zhou Y., Lai C.Y., Ren C. (2025). From open air to air-tight: Analyzing the ventilation overhaul in Hong Kong's wet markets and its implications. In \emph{Multiphysics and Multiscale Building Physics} (IBPC 2024), LNCE \textbf{555}, pp.\ 55--60. Springer; online December 2024. \doi{10.1007/978-981-97-8317-5_9}.}
\pub{C06}{\guo{}, Coleman J., Gullapalli I. (2025). Accelerating NZEB design optimization through LLM-based standardization and compliance checking. In \emph{Multiphysics and Multiscale Building Physics} (IBPC 2024), LNCE \textbf{553}, pp.\ 585--591. Springer; online December 2024. \doi{10.1007/978-981-97-8309-0_79}.}
\pub{C07}{\guo{}, Meggers F. (2019). Charging and discharging a coaxial borehole heat exchanger as a battery. \emph{Proceedings of Building Simulation 2019}, pp.\ 1749--1754. \doi{10.26868/25222708.2019.211375}.}
\pub{C08}{\guo{}, Teitelbaum E., Meggers F. (2019). Humidifying without adding humidity: Psychrometric shifts in humidity from air temperature setbacks enabled by radiant heating or cooling. \emph{Proceedings of Building Simulation 2019}, pp.\ 2562--2568.}
\pub{C09}{\guo{} (2018). Analytical evaluation of thermal comfort implications from an expanded range of air temperatures and velocities. \emph{15th International Conference on Indoor Air Quality and Climate (Indoor Air 2018)}.}


\pub{C10}{Coleman J., Teitelbaum E., \guo{}, Read J., Meggers F. (2017). Examining Architectural Air and Temperature with Novel Sensing Techniques. \emph{Energy Procedia}, \textbf{122}, 1135--1140. \doi{10.1016/j.egypro.2017.07.442}.}

\pub{C11}{\guo{}, Meggers F. (2017). Geothermal distribution network modeled as Heat Exchanger Network to be optimized with Mixed Integer Nonlinear Programming. \emph{Energy Procedia}, \textbf{122}, 1105--1110. \doi{10.1016/j.egypro.2017.07.466}.}

\pub{C12}{Meggers F., Teitelbaum E., Pantelic J., Rysanek A., \guo{} (2017). Liquid Desiccant Latent Load Handling Simulation for Building HVAC Applications with a DOAS Module. \emph{Proceedings of Building Simulation 2017}. \doi{10.26868/25222708.2017.616}.}

\pub{C13}{Teitelbaum E., \guo{}, Read J., Meggers F. (2017). Mapping Comfort with the SMART (Spherical Motion Average Radiant Temperature) Sensor. \emph{Proceedings of Building Simulation 2017}, \textbf{15}. \doi{10.26868/25222708.2017.644}.}

\pub{C14}{\guo{}, Meggers F. (2015). Impact of Control Availability on Perceived Comfort. \emph{Energy Procedia}, \textbf{78}, 1671--1677. \doi{10.1016/j.egypro.2015.11.254}.}

\pub{C15}{Calabrò E., Meggers F., Teitelbaum E., \guo{}, Gmachl C., Penello G.M. (2015). Thermoheliodome Testing: Evaluation Methods for Testing Directed Radiant Heat Reflection. \emph{Energy Procedia}, \textbf{78}, 1762--1768. \doi{10.1016/j.egypro.2015.11.298}.}

\pub{C16}{Meggers F., Aschwanden G., Teitelbaum E., \guo{}, Bruelisauer M. (2015). Urban Cooling Potential: System Losses from Microclimates. \emph{Energy Procedia}, \textbf{78}, 3072--3077. \doi{10.1016/j.egypro.2015.11.759}.}

\end{publications}

\subsection*{Manuscripts under review}
\begin{publications}
\pub{M01}{\guo{}, He K., Xu Y., Shi Z., Aviv D. (2026). Probabilistic thermal comfort for energy-efficient building control: A risk-aware framework. \emph{Energy Conversion and Management}. Under review.}
\end{publications}

\section{Patent applications}
\cvrow{2022}{US 2022/0300860 A1 / Co-inventor}
  {Methods and systems for generating predictions based on time series data using an ensemble model.
   Patent application published 2022.}
\cvrow{2022}{US Provisional 63/300,130 / Co-inventor}
  {Wearable metabolics: Human heat dissipation sensor via non-intrusive core and skin temperature
   measurement with custom-fitting wearables. Filed 2022.}
